\input{../tex-inputs/latex-leseansicht-vorspann}

\section[Arthur und Olga Schnitzler an Gustav Schwarzkopf, 4. 7. 1906]{L04385 Arthur und Olga Schnitzler an Gustav Schwarzkopf, 4. 7. 1906}
\nopagebreak\mylabel{L04385v}
\rehead{ }\normalsize\beginnumbering\briefempfaengerindex{Schwarzkopf, Gustav@\textsc{Schwarzkopf, Gustav}!zzzSchnitzler, Olga@\emph{von Olga Schnitzler}!1906-07-041@{4. 7. 1906}|(be}\briefempfaengerindex{Schwarzkopf, Gustav@\textsc{Schwarzkopf, Gustav}!zzzSchnitzler, Arthur@\emph{von Arthur Schnitzler}!1906-07-041@{4. 7. 1906}|(be}
\toendnotes[C]{\smallbreak\pagebreak[2]}
\correspDesc{Versand  durch Arthur Schnitzler, Olga Schnitzler am 4. 7. 1906 in Marienlyst
\newline{}Übermittlung  am 4. 7. 1906 in Marienlyst
\newline{}Erhalt  durch Gustav Schwarzkopf im Zeitraum [5. 7. 1906 – 9. 7. 1906?] in Tiefer Graben 23}\toendnotes[C]{\smallbreak}
\Standort{DLA, A:Schnitzler, HS.1985.1.1897.}
\physDesc{Bildpostkarte, 240 Zeichen
\newline{}Handschrift Arthur Schnitzler: Bleistift, deutsche Kurrent
\newline{}Handschrift Olga Schnitzler: Bleistift
\newline{}Versand: Stempel: »\nobreak{}4. 7. 0\textcolor{gray}{6}, 2–6S\nobreak{}«.  }\toendnotes[C]{\smallbreak}\pstart{}{\pb}Austria\oindex{Österreich@\textbf{Österreich}|pw}\pend{}\pstart{}Hr Gustav Schwarzkopf\pend{}\pstart{}Wien I\oindex{I., Innere Stadt@\textbf{I., Innere Stadt}|pw}\pend{}\pstart{}Tiefer Graben 23\oindex{Wien@\textbf{Wien}!I., Innere Stadt@\textbf{I., Innere Stadt}!Tiefer Graben 23@\textbf{Tiefer Graben 23}|pw}.\pend{}{\bigskip}
\pstart
           \centering{}{\pb}\textcolor{gray}{\textbf{Hamlets Grav\oindex{Hamlets Grav@\textbf{Hamlets Grav}|pw}.}}\pend
           \vspace{1em}
\pstart
           {\pb}Wo Roſenkranz\pwindex{Shakespeare, William 23.\,4.\,1564? Stratford-upon-Avon – 3.\,5.\,1616 ebd.@\textsc{Shakespeare, William} (23.\,4.\,1564? Stratford-upon-Avon – 3.\,5.\,1616 ebd.)!Hamlet@\strich\emph{Hamlet}|pwv} und Güldenſtern\pwindex{Shakespeare, William 23.\,4.\,1564? Stratford-upon-Avon – 3.\,5.\,1616 ebd.@\textsc{Shakespeare, William} (23.\,4.\,1564? Stratford-upon-Avon – 3.\,5.\,1616 ebd.)!Hamlet@\strich\emph{Hamlet}|pwv} liegen, weiſs man nicht ſicher\pend
           
\pstart
           {\pb}\textsc{Marienlyst\oindex{Marienlyst@\textbf{Marienlyst}|pw}}, 4/7 906\pend
           \vspace{0.5em}
\pstart
           herzlichen Gruß. Es iſt wunderſchön hier u wir dürften noch einige Zeit bleiben.\pend
           
\pstart
           Grüße an Ihren Bruder\pwindex{Schwarzkopf, Max 12.\,6.\,1857 Wien – 14.\,4.\,1928 ebd.@\textsc{Schwarzkopf, Max} (12.\,6.\,1857 Wien – 14.\,4.\,1928 ebd.)|pwv}.\pend
           
\pstart
           Ihr{\\[\baselineskip]}\spacefill\mbox{Arthur}\pend
           \leftskip=0em{}\selectlanguage{ngerman}\vspace{1em}\pstart \spacefill\mbox{{[}hs. O. Schnitzler:{]} O. S.}\pend{}\selectlanguage{ngerman}\endnumbering\briefempfaengerindex{Schwarzkopf, Gustav@\textsc{Schwarzkopf, Gustav}!zzzSchnitzler, Olga@\emph{von Olga Schnitzler}!1906-07-041@{4. 7. 1906}|)be}\briefempfaengerindex{Schwarzkopf, Gustav@\textsc{Schwarzkopf, Gustav}!zzzSchnitzler, Arthur@\emph{von Arthur Schnitzler}!1906-07-041@{4. 7. 1906}|)be}\mylabel{L04385h}
\begin{anhang}
\end{anhang}\newcommand{\dateiname}{L04385}\newcommand{\titel}{Arthur und Olga Schnitzler an Gustav Schwarzkopf, 4. 7. 1906}\newcommand{\editorInnen}{Herausgegeben von Jahnke, SelmaMüller, Martin Anton}\input{../tex-inputs/latex-leseansicht-abspann}
